% Created 2020-12-13 Sun 20:24
% Intended LaTeX compiler: pdflatex
\documentclass[11pt]{article}
\usepackage[utf8]{inputenc}
\usepackage[T1]{fontenc}
\usepackage{graphicx}
\usepackage{grffile}
\usepackage{longtable}
\usepackage{wrapfig}
\usepackage{rotating}
\usepackage[normalem]{ulem}
\usepackage{amsmath}
\usepackage{textcomp}
\usepackage{amssymb}
\usepackage{capt-of}
\usepackage{hyperref}
\usepackage{color}
\usepackage{listings}
\usepackage[english]{babel}
\author{simmi}
\date{\today}
\title{}
\hypersetup{
 pdfauthor={simmi},
 pdftitle={},
 pdfkeywords={},
 pdfsubject={},
 pdfcreator={Emacs 27.1 (Org mode 9.3)}, 
 pdflang={Germanb}}
\begin{document}

\titlepage \tableofcontents \clearpage\renewcommand\thesection{}
\section{Vorwort}
\label{sec:org7b62de3}
\cb{A}lle Jahre wieder! Ach ja, Weihnachten, da war doch was! Wenn du
das hier gerade liest und dich sprichwörtilch oder wortwörtlich am
Kopf kratzt und dich fragst, \emph{was soll das hier jetzt?}, dann habe ich
es wohl versäumt dich in den vergangenen Jahren mit einer meiner
Weihnachtsgeschichten zu beehren. Aber nun ist es endlich soweit! Ganz
toll, aber was kannst du dir jetzt darunter vorstellen, was erwartet
dich? Auch wenn das für die, die bereits in den vergangenen Jahren in
das \emph{Vergnügen} gekommen sind, ein wenig redundant ist werde ich nun
erstmal erklären was es sich mit diesem Schriftstück auf sich hat. Für
alle, die keine Weihnachtsgeschichten-Erstis sind, ihr könnt gerne zum
letzten Absatz des Vorworts springen.

\noindent kkkk

\setcounter{section}{0}
\renewcommand\thesection{Kapitel \arabic{section}:}
\end{document}